\documentclass[10pt,a4paper]{amsart}
\usepackage[margin=19mm]{geometry}
\usepackage{amsmath,amssymb,mathtools}
\usepackage[scr=rsfs,cal=euler]{mathalfa}
\usepackage{xfrac}
\usepackage[unicode,hidelinks,bookmarks=false]{hyperref}
\newcommand{\ie}{i.e.\ }
\newcommand{\cf}{cf.\ }
\newcommand{\resp}{respectively\ }
\newcommand{\Subsection}{\subsection}
\providecommand{\nobreakdash}{\nobreak\hbox{-}}
\setlength{\emergencystretch}{2em}
\multlinegap0pt
\usepackage{amssymb}
\newtheorem{proposition}{Proposition}
\DeclareMathOperator{\sgn}{\mathrm{sgn}}
\title{Multicrossing complex of knot and secant classes}
\author{Igor Nikonov}
\date{Source: arXiv:2609.00285v1 (CC BY 4.0)}
\begin{document}
\maketitle

\noindent\emph{Source and licence.} Multicrossing complex of knot and secant classes. Version arXiv:2609.00285v1 (CC BY 4.0), distributed by arXiv under the Creative Commons Attribution 4.0 International licence, http://creativecommons.org/licenses/by/4.0/, as recorded in the arXiv licence metadata for this exact version and shown on the abstract page https://arxiv.org/abs/2609.00285v1. Full licence notice: \texttt{licenses/arxiv-2609.00285-CC-BY-4.0.txt}.\par\smallskip

\noindent\textit{Editorial scope.} Counted unit: the article's proposition prop:4secant\_sign\_geometry (source0.tex lines 351--356). The article's complete original proof of it is delivered with it (source0.tex lines 358--427, one \texttt{proof} environment of 70 lines). No mathematics is written, repaired or re-derived: the statement and the proof are the article's own lines, in the article's own notation, and no retained line is substituted or re-flowed.  No further source-local context is needed: every cross-reference of every retained line resolves inside the two delivered spans.  One declared adaptation: exactly 1 ancillary strings inside the retained spans are omitted (image-inclusion commands and, where the source repeats a label, the repeated label definitions) and no other line differs from the source; the figure environments, with their placement, captions and labels, are kept, and the package records the omitted strings in fidelity receipts bound to the affected bodies.  No retained line contains a citation, so the wrapper carries no bibliography and no bibliography entry.  The retained lines contain the article's own diagram code, which the wrapper typesets with the standard tikz packages; no image file is delivered and the package declares no asset dependency.  Editorial formatting only, in addition: an AMS \texttt{amsart} preamble that restates the article's own declarations and macros used by the retained lines, namely the environments \texttt{proposition} on separate counters, together with the standard packages those lines need.  The retained environments print in this document as Proposition 1 in order of appearance, whereas the article prints the same items with the numbering of its own full document.  The complete native source of the article as distributed in the arXiv e-print for this exact version is bundled unchanged as \texttt{sources/209013/source0.tex}.
\begin{proposition}[Geometric meaning of the quadrisecant sign]\label{prop:4secant_sign_geometry}
\begin{multline*}
\sgn\Delta_4(\alpha_1,\alpha_2,\alpha_3,\alpha_4)=\\
or(\alpha_{12\underline{3}},\alpha_{12\underline{4}})\cdot or(\alpha_1,\alpha_2)\cdot\sgn[(z_1-z_3)(z_2-z_3)(z_1-z_4)(z_2-z_4)].
\end{multline*}
\end{proposition}

\begin{proof}
    First, we consider the case where the plane $\Pi$ is parallel to $\dot\alpha_1$ and $\dot\alpha_2$. Since the left term and the right term of the equation are infinitesimal, we can linearize the curves: $\alpha_1(u)=\alpha_1+\dot\alpha_1u$, $\alpha_2(v)=\alpha_2+\dot\alpha_2v$.

    Consider the coordinate system with the center $O$ and the basis $\dot\alpha_1$, $\dot\alpha_2$ and $\kappa$. The bundle of secants on the curves $\alpha_1(u)$ and $\alpha_2(v)$ projects a point $(x,y,z)$ to a point $(x_0,y_0,0)\in\Pi$ (Fig.~\ref{fig:secant_sheaf}).

\begin{figure}
    \centering
    
    \caption{2-secant projection}
    \label{fig:secant_sheaf}
\end{figure}

The secant intersects $\alpha_1$ at a point $(u,0,z_1)$ and intersects $\alpha_2$ at a point $(0,v,z_2)$. Then
\begin{equation}\label{eq:xy_to_uv}
u=\frac{z_1-z_2}{z-z_2}x,\quad v=\frac{z_2-z_1}{z-z_1}y
\end{equation}

and
\[
x_0=-\frac{z_2}{z_1-z_2}u=-\frac{z_2}{z-z_2}x,\quad y_0=-\frac{z_2}{z_2-z_1}v=-\frac{z_1}{z-z_1}y.
\]

Let $\alpha(t)$ be a curve such that $\alpha(0)$ lies on the quadrisecant $l$, i.e. $\alpha(0)=(0,0,z(0))$. Then $\alpha(t)=x(t)\dot\alpha_1+y(t)\dot\alpha_2+z(t)\kappa$ where

\begin{equation}\label{eq:alpha_to_xyz}
x(t)=\frac{(\kappa,\alpha(t),\dot\alpha_2)}{(\dot\alpha_1,\dot\alpha_2,\kappa)},\quad
y(t)=-\frac{(\kappa,\alpha(t),\dot\alpha_1)}{(\dot\alpha_1,\dot\alpha_2,\kappa)},\quad
z(t)=\frac{(\dot\alpha_1,\dot\alpha_2,\alpha(t))}{(\dot\alpha_1,\dot\alpha_2,\kappa)}.
\end{equation}

Then the derivations at $t=0$ are equal
\begin{equation}\label{eq:x0y0_derivation}
\begin{gathered}
    \dot x_0=-\frac{z_2}{z-z_2}\dot x=-\frac{z_2}{z-z_2}\cdot\frac{(\kappa,\dot\alpha,\dot\alpha_2)}{(\dot\alpha_1,\dot\alpha_2,\kappa)},\\
    \dot y_0=-\frac{z_1}{z-z_1}\dot y=\frac{z_1}{z-z_1}\cdot\frac{(\kappa,\dot\alpha,\dot\alpha_1)}{(\dot\alpha_1,\dot\alpha_2,\kappa)}.
\end{gathered}
\end{equation}

Applying~\eqref{eq:x0y0_derivation} to the curves $\alpha_3(t)$ and $\alpha_4(t)$, we see that $or(\alpha_{12\underline{3}},\alpha_{12\underline{4}})$ is the sign of the expression
\[
or(\alpha_1,\alpha_2)\cdot\begin{vmatrix}
                                 x_{03} & x_{04}\\ y_{03} & y_{04}
                             \end{vmatrix}
\]
where
\begin{multline*}
\begin{vmatrix}
x_{03} & x_{04}\\ y_{03} & y_{04}
\end{vmatrix}=
\begin{vmatrix}
   -\frac{z_2}{z_3-z_2}\cdot\frac{(\kappa,\dot\alpha_3,\dot\alpha_2)}{(\dot\alpha_1,\dot\alpha_2,\kappa)} & -\frac{z_2}{z_4-z_2}\cdot\frac{(\kappa,\dot\alpha_4,\dot\alpha_2)}{(\dot\alpha_1,\dot\alpha_2,\kappa)} \\
  \frac{z_1}{z_3-z_1}\cdot\frac{(\kappa,\dot\alpha_3,\dot\alpha_1)}{(\dot\alpha_1,\dot\alpha_2,\kappa)} & \frac{z_1}{z_4-z_1}\cdot\frac{(\kappa,\dot\alpha_4,\dot\alpha_1)}{(\dot\alpha_1,\dot\alpha_2,\kappa)}
\end{vmatrix}=\\
-\frac{z_1z_2}{(\dot\alpha_1,\dot\alpha_2,\kappa)^2(z_1-z_3)(z_2-z_3)(z_1-z_4)(z_2-z_4)}\cdot\\
\cdot\begin{vmatrix}
   (z_1-z_3)(\kappa,\dot\alpha_2,\dot\alpha_3) & (z_1-z_4)(\kappa,\dot\alpha_2,\dot\alpha_4) \\
  (z_2-z_3)(\kappa,\dot\alpha_1,\dot\alpha_3) & (z_2-z_4)(\kappa,\dot\alpha_1,\dot\alpha_4)
\end{vmatrix}=\\
\frac{z_1z_2}{(\dot\alpha_1,\dot\alpha_2,\kappa)^2(z_1-z_3)(z_2-z_3)(z_1-z_4)(z_2-z_4)}\Delta_4(\alpha_1,\alpha_2,\alpha_3,\alpha_4).
\end{multline*}
Since $z_1<0$,  $z_2<0$ and $(\dot\alpha_1,\dot\alpha_2,\kappa)^2>0$,
\begin{multline*}
or(\alpha_{12\underline{3}},\alpha_{12\underline{4}})=\\
or(\alpha_1,\alpha_2)\cdot\sgn\Delta_4(\alpha_1,\alpha_2,\alpha_3,\alpha_4)
\cdot\sgn[(z_1-z_3)(z_2-z_3)(z_1-z_4)(z_2-z_4)].
\end{multline*}
that implies the statement for $\Pi$ parallel to $\dot\alpha_1$ and $\dot\alpha_2$.

The trisecants on $\alpha_1(t)$, $\alpha_2(t)$, $\alpha_3(t)$, and those on $\alpha_1(t)$, $\alpha_2(t)$, $\alpha_4(t)$ form two surfaces which intersect in the line $l$. Then $or(\alpha_{12\underline{3}},\alpha_{12\underline{4}})$ is the sign of the dihedral angle between the surfaces and does not depend on the choice of the transversal plane $\Pi$.
\end{proof}

\end{document}
